% =====================================================================
%  29.3  Functional graphs: cycles with trees on top.
%  Exports: sec:c31-functional  def:c31-func
%  thm:c31-rhostruct  lem:c31-peel  lst:c31-peel  fig:c31-func
%  WORKED section + a solved lab.  In-section labs: 29.1 CF 118E, 29.2 CF 962F,
%  and here CF 702E "Analysis of Pathes in Functional Graph" (2100, L5) -- a
%  functional graph f on 0..n-1 with arc weights w and a fixed step count k
%  (k up to 1e10); for every start vertex output the SUM and MINIMUM arc weight
%  over the k-step walk.  Solved by binary lifting on f that carries a
%  (sum, min) monoid (the hardeningbox's own suggestion), O(n log k).  Live-
%  verified against the CF page (2s, 512 MB, n<=1e5, k<=1e10, 0<=w<=1e8); all
%  THREE official samples replayed (7 3 -> 10 1/8 1/7 1/10 2/8 2/7 1/9 3;
%  4 4 -> 0 0/4 1/8 2/12 3; 5 3 -> 7 1/17 1/19 2/21 3/8 1).  Promoted from this
%  chapter's mixed set, so its URL now lives here ONLY (no duplicate).
%  A functional graph (a set f: V -> V, out-degree exactly 1) has each
%  weak component shaped like a rho: one directed cycle with rooted in-trees
%  feeding it.  Detection by in-degree peeling (a queue, Ch.17's family) leaves
%  the cycle nodes; the stripped trees are ordinary rooted trees for Ch.13 DP.
%  rho(v)/mu(v) computed by DP.  Running object: f on 1..8.
% =====================================================================
\section{Functional Graphs: Cycles with Trees on Top}\label{sec:c31-functional}

\begin{previouslybox}[The graph had no directions; now every vertex has one arrow]
Sections 29.1--29.2 undid undirected redundancy by cutting edges or vertices.  A
\term{functional graph} is the opposite extreme of structure: every vertex has
\emph{exactly one} outgoing arrow, \(v\mapsto f(v)\), and the whole "graph" is just a
function.  Out-degree one is restrictive enough that the components are completely
classifiable: each is a single directed cycle with ordinary rooted trees growing
\emph{into} it -- a shape nicknamed the \term{rho} (\(\rho\)), because drawing a
tail that curls into a loop looks like the letter.  Once you find the cycle, the
trees are the trees of Chapters 12--13, and the arrows give you a canonical root
for free.
\end{previouslybox}

\begin{topicintro}
Every weak component of a functional graph has exactly one directed cycle
(\(\thmref{c31-rhostruct}\)): follow arrows from any vertex and, on a finite set, you
must eventually repeat -- the first repeat is the cycle.  The cleanest way to
\emph{extract} all cycles at once is \term{in-degree peeling}: push every vertex of
in-degree \(0\) to a queue, delete it, decrement its out-neighbor's in-degree, and
repeat; what survives every peel is precisely the union of cycles
(\(\lemref{c31-peel}\)) because every non-cycle vertex is a leaf of some in-tree and
gets peeled.  The remaining trees root on the cycle and are the natural domain of
the Ch.~13 bottom-up fold.  Distances to the cycle \(\rho(v)\) and the cycle length
\(\mu(v)\) fall out of a short DP (or binary lifting when you also get queries).

\textbf{Prerequisites.} Ch.~13 (\chref{c10}) for the fold on each in-tree; Ch.~17
(\chref{c28})/topological order for the peel queue.  Binary lifting (Ch.~16,
\chref{c11}) is optional for the query variant.

\textbf{After this section you can}
\begin{itemize}[nosep]
  \item recognise a functional graph in a statement (``each \(x\) has one \(f(x)\)'', ``repeatedly apply''); \hfill\lv{3}
  \item peel in-degrees to isolate every directed cycle in \(\bigO n\); \hfill\lv{4}
  \item compute \(\rho(v)\) and \(\mu(v)\) for all \(v\) with one DP; \hfill\lv{4}
  \item run Ch.~13 DP on each in-tree hanging off the cycle. \hfill\lv{4}
\end{itemize}
\end{topicintro}

% ---------------------------------------------------------------
\subsection{The rho decomposition}
% ---------------------------------------------------------------

\begin{definition}{Functional graph and its rho pieces}{c31-func}
A \term{functional graph} is a directed graph on \(V\) where each \(v\) has out-degree
exactly one, i.e.\ the edges of a total function \(f\colon V\to V\).  Its
\term{weak components} (ignore arrow directions) each contain exactly one directed
cycle; deleting the cycle's edges leaves rooted in-trees hanging on the cycle
vertices.  Writing \(\rho(v)\) for the number of arrows from \(v\) until the first
cycle vertex reached and \(\mu(v)\) for the length of that cycle, vertices with
\(\rho(v)=0\) are the cycle itself (\(\defref{c31-func}\)).
\end{definition}

\begin{theorem}{One cycle per component}{c31-rhostruct}
A finite functional graph's weak component with \(k\) vertices contains exactly one
directed cycle, of some length \(\mu\ge1\), and \(k-\mu\) in-tree vertices feeding it.
\end{theorem}

\begin{proofbox}
Start anywhere and iterate \(f\); the finite vertex set forces a repeat, and the
first repeated vertex closes a directed cycle (walk in, loop).  Every vertex reaches
a cycle (iterate until it does), and it reaches \emph{one}: if a vertex's orbit
entered two cycles they'd be mutually reachable, hence one cycle.  Contract each
cycle to a supernode; every remaining vertex has out-degree one toward a
strictly-closer-to-cycle vertex, which forbids a second cycle and makes the
pre-cycle vertices an in-tree (edges point toward the root). \(\Box\)
\end{proofbox}

\begin{lemma}{In-degree peeling leaves exactly the cycles}{c31-peel}
Repeatedly deleting a vertex of (current) in-degree \(0\) and its outgoing edge
terminates with the vertex set being exactly the union of directed cycles.
\end{lemma}

\begin{proofbox}
A vertex of in-degree \(0\) is a leaf of an in-tree (nothing points to it), hence not
on a cycle, so deleting it is safe and it cannot be the reason another cycle vertex
survives.  Conversely every in-tree has a leaf (a vertex of maximum depth), so as long
as any non-cycle vertex remains, some in-degree-\(0\) vertex exists to peel; peeling
shrinks the in-tree and never creates a new cycle vertex.  When peeling halts, every
remaining vertex has in-degree \(\ge1\) within the remainder -- but out-degree is still
one, so the remainder is a disjoint union of directed cycles. \(\Box\)
\end{proofbox}

\begin{lstlisting}[caption={Functional-graph decomposition: peel in-degrees to find cycle nodes, then a DFS/DP gives rho(v) tail length and mu(v) cycle length per component. ASCII only, identifiers without underscores}, label={lst:c31-peel}]
#include <bits/stdc++.h>
using namespace std;
const int MAXN=200005;
int n, f[MAXN], indeg[MAXN], mu[MAXN], rho[MAXN], cycLen[MAXN], onCycle[MAXN];
void peel(){
  queue<int> q;
  for(int v=1;v<=n;v++) if(indeg[v]==0) q.push(v);
  while(!q.empty()){
    int v=q.front(); q.pop(); onCycle[v]=0;          // peeled => not on a cycle
    int u=f[v]; if(--indeg[u]==0) q.push(u);          // its only out-edge removed
  }
  // whatever still has indeg>0 is on a cycle
  for(int v=1;v<=n;v++) if(indeg[v]>0) onCycle[v]=1;
}
void labelCycles(){
  for(int v=1;v<=n;v++) mu[v]=rho[v]=0;
  for(int v=1;v<=n;v++) if(onCycle[v] && mu[v]==0){   // walk one fresh cycle
    int cur=v, len=0; while(!mu[cur]){ mu[cur]=1; cur=f[cur]; len++; }  // mark pass
    cur=v; for(int i=0;i<len;i++){ mu[cur]=len; rho[cur]=0; cur=f[cur]; } // real length
  }
}
void dfsRho(int v,int p,int d){                        // v is an in-tree node, p parent
  rho[v]=d; mu[v]=mu[p];                                // inherit the cycle it feeds
  for(int u=1;u<=n;u++) if(f[u]==v && !onCycle[u]) dfsRho(u,v,d+1);
}
int main(){
  scanf("%d",&n);
  for(int v=1;v<=n;v++){ scanf("%d",&f[v]); indeg[f[v]]++; }
  memset(onCycle,1,sizeof(onCycle)); peel(); labelCycles();
  for(int c=1;c<=n;c++) if(onCycle[c]) for(int u=1;u<=n;u++) if(f[u]==c && !onCycle[u]) dfsRho(u,c,1);
  for(int v=1;v<=n;v++) printf("%d rho=%d mu=%d\n", v, rho[v], mu[v]);
}
\end{lstlisting}

\begin{figure}[H]
\centering
\begin{tikzpicture}[diagram, scale=0.95, transform shape,
   cy/.style={circle, draw=warningred, fill=warningred!12, line width=1pt,
              minimum size=6.4mm, font=\scriptsize\bfseries, inner sep=1pt},
   tr/.style={circle, draw=primaryblue!70, fill=primaryblue!8, line width=0.7pt,
              minimum size=6.4mm, font=\scriptsize\bfseries, inner sep=1pt},
   every edge/.style={draw, -Latex, line width=0.9pt}]
  \node[cy] (v1) at (0,0) {1}; \node[cy] (v2) at (1.2,0.9) {2}; \node[cy] (v3) at (1.2,-0.9) {3};
  \node[tr] (v4) at (3.0,1.4) {4}; \node[tr] (v5) at (4.4,1.9) {5};
  \node[cy] (v6) at (0,-2.8) {6};
  \node[cy] (v7) at (3.2,-2.6) {7}; \node[cy] (v8) at (4.4,-2.6) {8};
  \draw (v1)--(v2) (v2)--(v3) (v3)--(v1);
  \draw (v4)--(v2) (v5)--(v4);
  \draw (v6)to[out=90,in=90,loop above] node[above,font=\tiny]{self} (v6);
  \draw (v7)--(v8) (v8)--(v7);
  \node[annot, font=\scriptsize, anchor=west] at (-1.6,0.2) {component 1: cycle 1-2-3, tails 4->2, 5->4};
  \node[annot, font=\scriptsize] at (0.4,-3.9) {6: mu=1};  \node[annot, font=\scriptsize] at (3.9,-3.6) {7-8: mu=2};
\end{tikzpicture}
\caption{The running functional map \(f\) on \(\{1,\dots,8\}\):
\(1{\to}2,2{\to}3,3{\to}1,4{\to}2,5{\to}4,6{\to}6,7{\to}8,8{\to}7\).  Three weak
components -- a 3-cycle with an in-tree of depth \(2\) hanging at \(5\), a self-loop,
and a 2-cycle.  Red = cycle, blue = in-tree.  \(\rho,\mu\) are traced below.}
\label{fig:c31-func}
\end{figure}

\begin{insightbox}[Why out-degree one is a gift, not a constraint]
On a general directed graph, ``does the process terminate?'' and ``which cycle do I
land on?'' are hard.  Out-degree one collapses every component to a rho, so the
eventual behaviour of iterating \(f\) is summarised by two integers per vertex,
\(\rho(v)\) and \(\mu(v)\), computable in total linear time.  Countless problems in
disguise are this: a permutation is a functional graph of only cycles (all
\(\rho=0\)); ``each student votes for one candidate'' is a functional graph on the
vote map; a pointer machine that always follows one field is a functional graph.
Spot the single outgoing arrow and the rho machinery applies verbatim.
\end{insightbox}

\subsection{Worked trace: \(\rho\) and \(\mu\) on the running map}

\begin{tracebox}[Peeling, then labelling, of \(f\) from \(\figref{c31-func}\)]
\textbf{In-degrees.} \(\mathrm{indeg}=(1{:}1,\,2{:}2,\,3{:}1,\,4{:}1,\,5{:}0,\,6{:}1,\,7{:}1,\,8{:}1)\).
Queue starts with the only in-degree-\(0\) vertex, \(5\).  Peel \(5\): delete \(5{\to}4\), so
\(\mathrm{indeg}[4]\) drops \(1\to0\); push \(4\).  Peel \(4\): delete \(4{\to}2\),
\(\mathrm{indeg}[2]\) drops \(2\to1\) (not \(0\), no push).  Queue empty.  Survivors with
\(\mathrm{indeg}>0\): \(\{1,2,3,6,7,8\}\) -- exactly the three cycle vertex sets
(\(\lemref{c31-peel}\)); \(4,5\) were correctly peeled off as in-tree.\\
\textbf{Cycle lengths.} Walk component \(\{1,2,3\}\): \(\mu(1)=\mu(2)=\mu(3)=3\),
\(\rho=0\).  Self-loop \(6\): \(\mu(6)=1\).  2-cycle \(7,8\): \(\mu(7)=\mu(8)=2\).\\
\textbf{In-tree tails.} From cycle vertex \(2\), the un-peeled child is \(4\):
\(\rho(4)=1,\mu(4)=3\).  From \(4\), child \(5\): \(\rho(5)=2,\mu(5)=3\).  Final table
\(\rho=(0,0,0,1,2,0,0,0)\), \(\mu=(3,3,3,3,3,1,2,2)\).  Iterating \(f\) from \(5\):
\(5{\to}4{\to}2{\to}3{\to}1{\to}2\ldots\) -- after \(\rho(5)=2\) steps you are on the
3-cycle, then you circle with period \(\mu(5)=3\). \(\Box\)
\end{tracebox}

\begin{complexitybox}[Functional decomposition]
In-degree peeling is \(\bigO n\) (each vertex and each of the \(n\) edges touched once).
Labelling cycles and the tail DP is another \(\bigO n\).  Total \(\bigO n\) time, \(\bigO n\)
space -- optimal, since the input is \(n\) arrows.  The in-tree DP (Ch.~13) rides on top;
the \code{for(int u: f[u]==v)} child scan in \(\coderef{c31-peel}\) should be a prebuilt
reverse-adjacency list to stay \(\bigO n\), not \(\bigO{n^2}\) -- see the warning below.
\end{complexitybox}

\begin{warningbox}[Three functional-graph errors]
\begin{enumerate}[nosep,leftmargin=1.6em]
  \item \textbf{Building children by scanning all of \(f\).} \(\coderef{c31-peel}\)'s
        \code{dfsRho} inner loop \code{if(f[u]==v)} over all \(u\) is \(\bigO n\) per
        node -- \(\bigO{n^2}\) total.  Build a reverse adjacency
        \code{radj[f[v]].push\_back(v)} once; then the in-tree DFS is \(\bigO n\).
  \item \textbf{Two-phase cycle marking.} The \code{labelCycles} pass first scribbles
        \(\mu{\neq}0\) as a visited flag, then rewrites the true length; forgetting the
        rewrite leaves every cycle length equal to \(1\).  Keep both passes.
  \item \textbf{Assuming one component.} A functional graph is a \emph{forest} of
        rhos; iterate the labeller over all unlabelled cycle vertices, not just vertex
        \(1\).  Symptom: only the first component gets \(\mu,\rho\).
\end{enumerate}
\end{warningbox}

\begin{hardeningbox}
\begin{itemize}[nosep,leftmargin=1.4em]
  \item \textbf{Queries without the peel.} If you must answer ``where is \(v\) after
        \(t\) arrows?'' online, use binary lifting on \(f\) (Ch.~16, \chref{c11}):
        \(\mathit{up}[j][v]=f^{\,2^j}(v)\) answers \(f^{\,t}(v)\) in
        \(\bigO{\log t}\).  The peel still gives \(\rho,\mu\) for cycle-membership
        queries in \(O(1)\) after \(\bigO n\).
  \item \textbf{Detecting a permutation.} \(f\) is a bijection iff every in-degree is
        exactly \(1\), i.e.\ no peel ever fires and every \(\rho(v)=0\); recognising
        this turns ``cycle decomposition'' problems (elementary symmetric sums of
        cycle lengths) into a straight read of \(\mu\) over cycle roots.
  \item \textbf{Synchronous iteration.} ``Do all vertices step at once'' problems need
        the whole map updated simultaneously (\(f\) applied to snapshot values); model
        it as the functional graph of the joint update and read cycle lengths to get
        the period.
\end{itemize}
\end{hardeningbox}

% ---------------------------------------------------------------
\subsection{Solved: CF 702E -- Analysis of Pathes in Functional Graph (2100, L5)}
% ---------------------------------------------------------------

\solvedproblem{c31-funcpath}{Analysis of Pathes in Functional Graph}{CF 702E}{https://codeforces.com/problemset/problem/702/E}{CF 2100 L5}{functional graph, binary lifting on \(f\) carrying a (sum, min) monoid, rho cycle structure}

\begin{problemstatement}[CF 702E, verified against the live CF problem page]
A functional graph on vertices \(0,\dots,n-1\) is given as arrays \(f_0,\dots,f_{n-1}\)
(\(0\le f_i<n\), the unique arc out of \(i\)) and arc weights \(w_0,\dots,w_{n-1}\)
(\(0\le w_i\le 10^8\), the weight of the arc \(i\to f_i\)).  A path length \(k\) is
also given.  For each start vertex \(i\), walk exactly \(k\) arcs following \(f\) and
output \(s_i\) = the \emph{sum} of the \(k\) arc weights on that walk and \(m_i\) =
the \emph{minimum} of those \(k\) arc weights.  Constraints: \(1\le n\le 10^5\),
\(1\le k\le 10^{10}\).  Time limit 2 seconds, memory 512 MB.  Samples:
\(n{=}7,k{=}3\), \(f=(1,2,3,4,3,2,6)\), \(w=(6,3,1,4,2,2,3)\) \(\to\) lines
\(10\ 1; 8\ 1; 7\ 1; 10\ 2; 8\ 2; 7\ 1; 9\ 3\); \(n{=}4,k{=}4\),
\(f=(0,1,2,3)\), \(w=(0,1,2,3)\) \(\to\) \(0\ 0; 4\ 1; 8\ 2; 12\ 3\);
\(n{=}5,k{=}3\), \(f=(1,2,3,4,0)\), \(w=(4,1,2,14,3)\) \(\to\) \(7\ 1; 17\ 1; 19\ 2;
21\ 3; 8\ 1\).
\end{problemstatement}

\begin{attemptgate}[Attempt it yourself first -- \lv{5}, give yourself 40 minutes]
\(k\) is up to \(10^{10}\): you cannot step through the walk per vertex.  The orbit
of \(i\) is a rho (\(\secref{c31-functional}\)), but the section's own hardeningbox
named a cleaner tool for exactly this "fixed power of \(f\)" query.  Which lift, and
what does each jump have to \emph{accumulate} so the sum and minimum fold correctly?
\peek{Binary lifting on \(f\) (Ch.~16, \(\chref{c11}\)): store \(f^{2^j}\) per vertex.
But a sum/min over a walk is not just a landing node -- each \(2^j\)-jump must also
carry the sum and the running minimum of the \(2^j\) arcs it covers, combined when
jumps are chained.}
\nudge{Define \(\mathit{to}[j][v]\) (node after \(2^j\) steps), \(\mathit{sum}[j][v]\)
(sum of those \(2^j\) arc weights) and \(\mathit{mn}[j][v]\) (their minimum).
\(\mathit{to}[0][v]=f(v)\), \(\mathit{sum}[0][v]=\mathit{mn}[0][v]=w_v\).  To double:
the \(2^j\) walk from \(v\) is the \(2^{j-1}\) walk from \(v\) then the \(2^{j-1}\)
walk from \(\mathit{to}[j{-}1][v]\) -- add the sums, take the min of the mins.}
\fullhint{Build the three tables for \(j=1..34\); then for vertex \(v\) sweep the set
bits of \(k\) from low to high, cur starting at \(v\), adding \(\mathit{sum}[j][\text{cur}]\)
and taking \(\min\) with \(\mathit{mn}[j][\text{cur}]\) and jumping
\(\text{cur}\leftarrow\mathit{to}[j][\text{cur}]\).  Output the accumulated \((S,M)\).
This is \(\bigO{n\log k}\) with no rho casework -- the sums reach \(k\cdot 10^8\sim 10^{18}\)
so use \texttt{long long}.  Full build and the three sample replays are below.}
\end{attemptgate}

\solstep{Understanding}
Every vertex walks a fixed \(k\)-step orbit; doing it literally is \(\bigO{nk}\),
about \(10^{15}\).  The rho structure (tail plus \(\mu\) periodicity) admits an
\(\bigO n\) prefix-sum solution, but it needs careful casework (where the \(k\)-window
crosses from tail into cycle, and the minimum over a \emph{partial} final lap).  The
lift view from the hardeningbox folds the whole orbit uniformly: precompute
power-of-two jumps, then compose the \(\approx 34\) bits of \(k\).  The monoid being
carried is \((\mathrm{sum},\mathrm{min})\) -- both associative, so chaining jumps is clean.

\solstep{Key Observations}
\begin{observation}
The \(2^j\)-step walk from \(v\) splits at its midpoint: land \(\mathit{mid}=\mathit{to}[j{-}1][v]\)
after \(2^{j-1}\) steps, then take another \(2^{j-1}\) from \(\mathit{mid}\).  Hence
\(\mathit{sum}[j][v]=\mathit{sum}[j{-}1][v]+\mathit{sum}[j{-}1][\mathit{mid}]\) and
\(\mathit{mn}[j][v]=\min(\mathit{mn}[j{-}1][v],\mathit{mn}[j{-}1][\mathit{mid}])\).
\end{observation}
\begin{observation}
Writing \(k\) in binary, the \(k\)-walk is the concatenation, in increasing bit
order, of the \(2^j\)-walks from the \emph{current} node.  Sum adds and min folds, so
sweeping set bits from low to high while advancing \(\text{cur}\) accumulates
\((S,M)\) exactly -- no dependence on where the rho's tail ends.
\end{observation}
\begin{observation}
\(k<2^{34}\), so \(35\) levels suffice.  The tables are \(35\cdot n\) entries
\(\approx 3.5\cdot 10^6\) -- about \(70\) MB with two \texttt{long long} arrays and
one \texttt{int}, comfortably inside 512 MB.
\end{observation}

\solstep{Think It With Me}
\begin{thinkbox}
\deadend{Walk \(k\) arcs per vertex}{\(\bigO{nk}\approx 10^{15}\) steps -- \emph{TLE}.
Only works for tiny \(k\).}
\deadend{Peel the rho and prefix-sum around each cycle}{Correct and \(\bigO n\), but the
\(k\)-window's split into tail, whole laps, and a \emph{partial} closing lap must be
handled for the MINIMUM separately -- three error-prone cases.}
\aha{Lift with a \((\mathrm{sum},\mathrm{min})\) monoid: \(\bigO{n\log k}\), one uniform
fold, no casework.  The associativity of sum and min is what makes the jump chaining legal.}
\end{thinkbox}

\solstep{The Approach Ladder}
\begin{center}\renewcommand{\arraystretch}{1.25}\setlength{\tabcolsep}{3pt}
\begin{tabular}{@{}c p{0.28\linewidth} p{0.30\linewidth} l c@{}}
\toprule
\textbf{\#} & \textbf{Approach} & \textbf{Idea} & \textbf{Time} & \textbf{Verdict}\\
\midrule
1 & \lv{2}\ step walk & follow \(f\) \(k\) times per vertex & \(\bigO{nk}\) & \tle\\
2 & \lv{4}\ rho + cycle prefix & peel, prefix sum/min on the ring & \(\bigO n\) & \textbf{WA-prone}\\
3 & \lv{5}\ \textbf{binary lift (sum,min)} & fold bits of \(k\) & \(\bigO{n\log k}\) & \ac\\
\bottomrule
\end{tabular}\end{center}
Rung 3 trades a small memory for the elimination of all tail/lap/partial-lap casework
-- the reliable contest choice, and precisely the lift the section's hardeningbox
promised.

\solstep{Solution}
Read \(f\) into \(\mathit{to}[0][\cdot]\) and \(w\) into \(\mathit{sum}[0][\cdot]=
\mathit{mn}[0][\cdot]\).  For \(j=1..34\) build the doubling tables.  For each vertex
\(v\) start \(S=0, M=+\infty, \text{cur}=v\) and for every set bit \(j\) of \(k\) (low to
high) do \(S{+}=\mathit{sum}[j][\text{cur}]\), \(M=\min(M,\mathit{mn}[j][\text{cur}])\)
and \(\text{cur}=\mathit{to}[j][\text{cur}]\).  Print \(S\ M\).

\solstep{Implementation}
\begin{lstlisting}[caption={CF 702E. Binary lifting on the functional graph carrying a (sum, min) monoid; per vertex fold the set bits of k (up to 2 to the 34) low-to-high. Sums reach 1e18 so long long; min identity is a huge value}, label={lst:c31-funcpath}]
#include <bits/stdc++.h>
using namespace std;
typedef long long ll;
const int MAXN = 100005, B = 35;
int to_[B][MAXN];
ll sumv[B][MAXN], mnv[B][MAXN];
int main(){
  int n; ll k; scanf("%d %lld", &n, &k);
  for(int v=0; v<n; v++) scanf("%d", &to_[0][v]);
  for(int v=0; v<n; v++){ ll w; scanf("%lld", &w); sumv[0][v]=w; mnv[0][v]=w; }
  for(int j=1; j<B; j++)
    for(int v=0; v<n; v++){
      int mid = to_[j-1][v];
      to_[j][v]   = to_[j-1][mid];
      sumv[j][v]  = sumv[j-1][v] + sumv[j-1][mid];
      mnv[j][v]   = min(mnv[j-1][v], mnv[j-1][mid]);
    }
  for(int v=0; v<n; v++){
    ll S=0, M=(ll)4e18; int cur=v;
    for(int j=0; j<B; j++) if((k>>j)&1LL){
      S += sumv[j][cur]; M = min(M, mnv[j][cur]); cur = to_[j][cur];
    }
    printf("%lld %lld\n", S, M);
  }
  return 0;
}
\end{lstlisting}

\solstep{Trace on the Official Samples}
\begin{tracebox}[Sample 3: \, 5-cycle \(f=(1,2,3,4,0)\), \(w=(4,1,2,14,3)\), \(k=3\), start \(v=0\)]
\(k=3\) has bits \(0,1\).  Bit 0 from cur\(=0\): \(S+=w_0=4\), \(M=4\), cur\(\to1\).  Bit 1
from cur\(=1\) (a 2-step jump \(1\to2\to3\)): \(\mathit{sum}[1][1]=w_1+w_2=1+2=3\),
\(\mathit{mn}[1][1]=\min(1,2)=1\).  \(S=4+3=7\), \(M=\min(4,1)=1\).  Output \(\mathbf{7\ 1}\).
The same fold from \(v=3\): bits 0,1 give \(w_3=14\) then \(w_4+w_0=3+4\), so
\(S=21,\ M=3\) \(\to\) \(\mathbf{21\ 3}\).
\end{tracebox}
\begin{tracebox}[Sample 2: \, all self-loops \(f=(0,1,2,3)\), \(w=(0,1,2,3)\), \(k=4\)]
\(k=4\) is a single bit 2.  From \(v\), the \(2^2\)-jump stays on \(v\) and covers
\(4\) copies of \(w_v\): \(S=4w_v,\ M=w_v\).  Vertices \(v=0..3\) give
\(0\ 0;\ 4\ 1;\ 8\ 2;\ 12\ 3\).
\end{tracebox}
All seven Sample-1 lines and the five Sample-3 lines reproduce identically.

\solstep{Complexity}
\begin{complexitybox}
Building the tables is \(\bigO{n\log k}\) (\(35n\) operations) and answering all
vertices is the same \(\bigO{n\log k}\).  For \(n=10^5\) that is \(\sim 3.5\cdot 10^6\),
comfortable in 2 seconds.  Memory is \(35n\) entries across an \texttt{int} and two
\texttt{long long} tables \(\approx 70\) MB, within 512 MB.  Sums reach
\(k\cdot\max w\sim 10^{10}\cdot 10^8=10^{18}\), so \(S\) and the table sums are
\texttt{long long}; the running \(M\) starts at \(+\infty\).
\end{complexitybox}

\begin{hardeningbox}
\begin{itemize}[nosep,leftmargin=1.4em]
\item \textbf{Lift \emph{once}, not per bit of value.}  Build \(\mathit{to}\),
\(\mathit{sum}\) and \(\mathit{mn}\) up to \(\lceil\log_2 k\rceil\); a common bug is sizing
the loop by \(n\)
instead of by \(k\) (\(k\) can exceed \(n\)).
\item \textbf{Sweep bits low-to-high while advancing cur.}  The jump for bit \(j\)
must start where the previous jump \emph{landed}; using the original \(v\) each time
mis-splits the orbit.
\item \(\min\) \textbf{identity and overflow.}  Initialise \(M\) above \(10^8\) (e.g.\ \(4\cdot
10^{18}\)); accumulate \(S\) in \texttt{long long} (\(10^{18}\) overflows \texttt{int}).
\item \textbf{0-indexed vertices.}  CF 702E uses \(0\le f_i<n\); do not \texttt{--} the labels
as you would in a 1-indexed problem.
\end{itemize}
\end{hardeningbox}

\begin{recallbox}[29.3 in one breath]
\emph{Out-degree one \(\Rightarrow\) each component is a rho: exactly one directed
cycle with in-trees feeding it (\(\thmref{c31-rhostruct}\)).  Peel in-degree-\(0\)
vertices to strip every tree and leave the cycles (\(\lemref{c31-peel}\)); inherit
\(\mu(v)\) down the tree and set \(\rho(v)\) by depth (\(\coderef{c31-peel}\)), all in
\(\bigO n\).  The trees are Ch.~13 territory.  Read the single outgoing arrow in the
statement as the tell, and prebuild reverse adjacency so the tail DFS stays linear.}
\selftest{
  \item Why must every weak component have a cycle? \emph{(Iterating \(f\) from any vertex on a finite set forces a repeat, and the first repeat closes a directed cycle \(\thmref{c31-rhostruct}\).)}
  \item What survives in-degree peeling and why? \emph{Exactly the cycle vertices: every non-cycle node is an in-tree leaf at some depth and gets peeled, and survivors all keep in-degree >=1 with out-degree 1, so they are cycles \(\lemref{c31-peel}\).)}
  \item Give \(\rho(5),\mu(5)\) for the running map. \emph{(\(\rho(5)=2\) steps to reach the cycle, \(\mu(5)=3\), the length of the 1-2-3 cycle \(\figref{c31-func}\).)}
  \item In CF 702E, why does bit order low-to-high with a moving \(\text{cur}\) fold both the sum and the minimum correctly? \emph{(Each stored jump carries its own (sum, min); sum and min are associative, so concatenating the bit-jumps in orbit order composes them exactly.)}
}
\end{recallbox}
